\section*{第 \ref{ch:g5:biodiversity} 章——生物多样性}

\begin{solution}{exo:g5:biodiversity:1}
一种特定的\omterm{def:g1:living-or-not:living}{生物}——
个体结构相同，彼此繁殖。
例如：家麻雀、普通雏菊、
赤狐、花园蜗牛。
\end{solution}

\begin{solution}{exo:g5:biodiversity:2}
指一个地方所拥有的生命多样性——
主要看其\omterm{def:g5:biodiversity:species}{物种}数目及健康状况。开花的草地
每平方米可容纳数十种动\omterm{def:g1:plants-around-us:plant}{植物}，\omterm{def:g5:biodiversity:biodiversity}{生物多样性}高；
铺装的庭院几乎一种也没有。
\end{solution}

\begin{solution}{exo:g5:biodiversity:3}
约二百万种。由于每次严肃的考察
仍会发现尚未命名的\omterm{def:g5:biodiversity:species}{物种}——尤以甲虫为甚——
因此目录之外必定还有数百万种。
\end{solution}

\begin{solution}{exo:g5:biodiversity:4}
一个\omterm{def:g2:where-animals-live:habitat}{栖息地}内\omterm{def:g5:biodiversity:species}{物种}众多（草地的一平方米）；
多个\omterm{def:g2:where-animals-live:habitat}{栖息地}毗邻（池塘、树篱、草地）；
同一\omterm{def:g5:biodiversity:species}{物种}内部存在差异（没有两只花园蜗牛的壳完全相同）。
\end{solution}

\begin{solution}{exo:g5:biodiversity:5}
在丰富的食物网中，若某种猎物减少，
捕食者可转而捕食其他猎物，每种角色都有替补；冲击因此被吸收。
在贫乏的食物网中，每个\omterm{def:g5:biodiversity:species}{物种}可选择的余地很少，
因此一个\omterm{def:g5:biodiversity:species}{物种}的消失会连带拖累其他\omterm{def:g5:biodiversity:species}{物种}。
\end{solution}

\begin{solution}{exo:g5:biodiversity:6}
蜜蜂和食蚜蝇：为\omterm{def:g1:plants-around-us:parts}{花朵}传粉。瓢虫：
捕食蚜虫。\omterm{prop:g3:sorting-living-things:groups}{鸟类}：巡视并捕食\omterm{ex:g2:animals-grow-up:butterfly}{毛毛虫}。
蚯蚓：使土壤保持疏松。
\end{solution}

\begin{solution}{exo:g5:biodiversity:7}
\omterm{def:g2:where-animals-live:habitat}{栖息地}遭破坏（池塘被填干）；\omterm{def:g2:caring-for-nature:environment}{环境}遭污染（垃圾和化学品）；
捕捞过度（过度捕鱼）。\omterm{prop:g5:biodiversity:loss}{灭绝}指彻底消失——
在所有地方、永远消失。
\end{solution}

\begin{solution}{exo:g5:biodiversity:8}
开放性答案。草地或种植绿化的方块
通常可记录到十种或更多\omterm{def:g5:biodiversity:species}{物种}；铺装的方块则接近零——
这一差异正是测量的结果。
\end{solution}

\begin{solution}{exo:g5:biodiversity:9}
每个\omterm{def:g5:biodiversity:species}{物种}都是一根线：它消失，会使每个邻近\omterm{def:g5:biodiversity:species}{物种}——
无论是吃它的、被它传粉的，还是被它清理的——
可选择的余地减少，并削弱食物网
抵御下一次冲击的能力。由于没有人知道
一个\omterm{def:g5:biodiversity:species}{物种}连着多少根线，因此无人能事先断定
某次\omterm{prop:g5:biodiversity:loss}{灭绝}无害；\omterm{prop:g5:biodiversity:loss}{灭绝}既然永久，
错误便再也无法挽回。
\end{solution}

\begin{solution}{exo:g5:biodiversity:10}
逐一保护鹳鸟，等于保护\omterm{prop:g3:sorting-living-things:groups}{鸟类}；
恢复湿润草地，则重建了鹳鸟所\emph{依赖}的
\omterm{def:g2:caring-for-nature:environment}{环境}——\omterm{def:g2:where-animals-live:habitat}{栖息地}会同时供养其中的所有居民，
这正是青蛙、蜻蜓和草地花卉随鹳鸟一同回归的原因。
没有\omterm{def:g2:where-animals-live:habitat}{栖息地}，就没有居民；\omterm{def:g2:where-animals-live:habitat}{栖息地}恢复，
居民便会回来。
\end{solution}

\begin{solution}{exo:g5:biodiversity:11}
被移除的不是装饰，而是一支劳动力：
树篱和花卉\omterm{def:g5:biodiversity:species}{物种}庇护着传粉者与食虫者，
它们的服务在不付费时从不计入任何账目。
自然的工作是看不见的记账——持续、
无偿、无人注意——因此其代价
只有在工人被赶走、工作以账单的形式
回来时，才被发现。
\end{solution}
